% \section{Related Work} \label{sec:related_work}

% Cohesive subgraph discovery is a fundamental problem in graph mining, with numerous models proposed to capture various notions of cohesiveness. 
% One of the most widely used models is $k$-core decomposition~\cite{seidman1983network,zhang2023efficient,peel_cd,core_res1,kcore_4,io-eff,k-core3,core_res3,HyperCD}, which identifies dense regions in a subgraph by recursively removing vertices with degrees less than $k$, 
% Another popular model is $k$-truss decomposition~\cite{cohen2008trusses,tian2023truss,xu2022truss,chen2020truss,wang2012truss,wang2022truss,akbas2017truss,huang2014truss}, which uncovers tightly connected regions in a subgraph by recursively removing edges with support less than $k$, where the support of an edge is defined as the number of triangles containing that edge. 
% % 
% Nuclear decomposition is first introduced by \cite{firstNucleus}, which generalizes the well-known $k$-core($1,2$) and $k$-truss($2,3$) decompositions to higher-order $(r,s)$-nucleus relationships.
% Other works on nucleus decomposition have been discussed in previous sections and will not be repeated here.

% \myadd{
% Beyond core/truss-style peeling, dense subgraph mining and quasi-clique discovery study (near-)cliques under density-based objectives or relaxed constraints~\cite{Lanciano2024DensestSurvey,Konar2024OQC}. 
% These formulations are complementary to $(r,s)$-nuclei: while dense-subgraph/quasi-clique methods target an objective-driven dense subgraph (or a set of near-cliques), $(r,s)$-nucleus decomposition is defined via higher-order clique containment and support and produces a global hierarchy across core values.

% A closely related direction is community search, which is query-dependent and aims to return a cohesive subgraph containing given query vertices~\cite{SozioGionisKDD2010,SunPVLDB2025GICS}. 
% Compared to decomposition, community search focuses on answering per-query results efficiently (often with additional user constraints such as size, connectivity, or interactivity), whereas our work studies the global $(r,s)$-nucleus decomposition and targets efficient peeling with dynamic index updates.

% Finally, clique enumeration and clique counting are fundamental primitives in graph mining and are frequently used as building blocks for higher-order cohesiveness measures. 
% Recent work has revisited scalable $k$-clique counting by combining exact counting with sampling-based techniques to handle large graphs~\cite{YeSIGMOD2022Lightning}. 
% }

% \section{Conclusion and future work} \label{sec:conclusion}
% In this paper, we revisited nucleus decomposition and presented the first \textit{counting-based} framework that eliminates explicit $s$-clique enumeration.
% At the core of our approach is the \textit{Clique Path Index} (\cpi), a compact auxiliary structure that organizes the $s$-clique search space and enables direct counting, efficient dynamic updates, and connectivity maintenance throughout decomposition.
% Extensive experiments show that our method significantly accelerates both nucleus decomposition and hierarchy construction, while scaling effectively to larger $s$ values.

% \myadd{
% Our counting-and-indexing viewpoint suggests several natural directions for future work.
% First, it is well-suited for parallelization since the computations are largely local and additive: \textsc{BuildIndex} and \textsc{InitialCounting} can be parallelized by partitioning pivot vertices (or seed $r$-cliques), because CPI path generation and local support contributions are independent and can be merged via a final reduction; during peeling, \textsc{UpdateIndex}/\textsc{UpdateSupport} can be parallelized by processing removed $r$-cliques in batches with thread-local buffers and aggregated (atomic) support updates.
% More broadly, the same principle is not tied to cliques: a \emph{clique} is a vertex set where every pair is connected~\cite{diestel2000graph}, and CPI serves to index how smaller patterns are contained in larger supporting patterns so that supports can be updated by counting.
% This abstraction extends to richer graph classes by replacing cliques with an appropriate ``complete group pattern'', e.g., bicliques for bipartite graphs (closely related to $(\alpha,\beta)$-core models)~\cite{tian2024abcore}, hyperedge-induced group patterns for hypergraphs~\cite{HyperCD}, and temporal $\Delta$-cliques for temporal graphs~\cite{viard2016deltaclique}.
% We leave optimized multi-threaded implementations and these extensions as future work.
% }

% \section{Related Work} \label{sec:related_work}

% Cohesive subgraph discovery is fundamental in graph mining. 
% Classic peeling-based models include $k$-core decomposition~\cite{seidman1983network,zhang2023efficient,peel_cd,core_res1,kcore_4,io-eff,k-core3,core_res3,HyperCD}, which recursively removes vertices with degree $<k$, and $k$-truss decomposition~\cite{cohen2008trusses,tian2023truss,xu2022truss,chen2020truss,wang2012truss,wang2022truss,akbas2017truss,huang2014truss}, which recursively removes edges whose support (the number of triangles containing the edge) is $<k$. 
% Nuclear decomposition was introduced by~\cite{firstNucleus}, generalizing $k$-core($1,2$) and $k$-truss($2,3$) to higher-order $(r,s)$-nucleus relationships; other nucleus-decomposition works have been discussed earlier and are not repeated.

% \myadd{
% Beyond core/truss-style peeling, dense subgraph mining and quasi-clique discovery consider (near-)cliques under density objectives or relaxed constraints~\cite{Lanciano2024DensestSurvey,Konar2024OQC}. 
% They are complementary to $(r,s)$-nuclei: dense-subgraph/quasi-clique methods target objective-driven dense subgraphs (or sets of near-cliques), while $(r,s)$-nucleus decomposition is defined via higher-order clique containment/support and yields a global hierarchy across core values.

% Community search is query-dependent, returning a cohesive subgraph containing given query vertices~\cite{SozioGionisKDD2010,SunPVLDB2025GICS}. 
% It emphasizes efficient per-query answers (often with constraints such as size, connectivity, or interactivity), whereas we study global $(r,s)$-nucleus decomposition and focus on efficient peeling with dynamic index updates.

% Finally, clique enumeration/counting are key primitives for higher-order cohesiveness measures; recent work revisits scalable $k$-clique counting by combining exact counting with sampling to handle large graphs~\cite{YeSIGMOD2022Lightning}.
% }

% \section{Conclusion and future work} \label{sec:conclusion}
% In this paper, we revisited nucleus decomposition and presented the first \textit{counting-based} framework that eliminates explicit $s$-clique enumeration.
% Our key structure is the \textit{Clique Path Index} (\cpi), which organizes the $s$-clique search space to enable direct counting, efficient dynamic updates, and connectivity maintenance during decomposition.
% Extensive experiments show substantial speedups for both decomposition and hierarchy construction, and strong scalability to larger $s$.

% \myadd{
% Our counting-and-indexing viewpoint also suggests natural future directions.
% First, it is amenable to parallelization since computations are largely local and additive: \textsc{BuildIndex} and \textsc{InitialCounting} can parallelize over pivot vertices (or seed $r$-cliques), where CPI path generation and local support contributions are independent and merged by a final reduction; during peeling, \textsc{UpdateIndex}/\textsc{UpdateSupport} can process removed $r$-cliques in batches using thread-local buffers with aggregated (atomic) support updates.

% More broadly, the principle is not tied to cliques. 
% A \emph{clique} is a vertex set with all pairs connected~\cite{diestel2000graph}, and CPI’s role is to index how smaller patterns are contained in larger supporting patterns so that supports are updated by counting. 
% This abstraction extends to richer graph classes by replacing cliques with suitable ``complete group patterns'': bicliques for bipartite graphs (closely related to $(\alpha,\beta)$-core models)~\cite{tian2024abcore}, hyperedge-induced group patterns for hypergraphs~\cite{HyperCD}, and temporal $\Delta$-cliques for temporal graphs~\cite{viard2016deltaclique}. 
% We leave optimized multi-threaded implementations and these extensions as future work.
% }



\section{Related Work} \label{sec:related_work}

Cohesive subgraph discovery is fundamental in graph mining. 
Classic peeling-based models include $k$-core decomposition~\cite{seidman1983network,zhang2023efficient,peel_cd,core_res1,kcore_4,io-eff,k-core3,core_res3,HyperCD}, which removes vertices with degree $<k$, and $k$-truss decomposition~\cite{cohen2008trusses,tian2023truss,xu2022truss,chen2020truss,wang2012truss,wang2022truss,akbas2017truss,huang2014truss}, which removes edges whose triangle support is $<k$. 
Nucleus decomposition was introduced by~\cite{firstNucleus}, generalizing $k$-core($1,2$) and $k$-truss($2,3$) to $(r,s)$-nuclei; other nucleus-decomposition works are discussed earlier and not repeated.
%
\myaddRC{
At the exact algorithm level, prior nucleus methods mainly differ in how they handle local clique structure during support updates.
%
Some methods maintain explicit clique incidence.
%
Other methods reconstruct affected local neighborhoods on demand.
%
Still, they share the same high-level peeling skeleton.
%
They also share a reconstructive support-maintenance primitive.
%
Our work differs at this point.
%
We still use the same exact decomposition paradigm.
%
However, we replace reconstructive maintenance by a two-layer design.
%
\cpi is the representation layer.
%
DPCM is the exact maintenance layer.
}
% 
\myaddRA{
Dense-subgraph mining and quasi-clique discovery target objective-driven dense regions or relaxed clique constraints~\cite{Lanciano2024DensestSurvey,Konar2024OQC}; they are complementary to $(r,s)$-nuclei, which are defined via higher-order clique containment/support and induce a global hierarchy across core values.
Community search is query-dependent, returning a cohesive subgraph containing given query vertices~\cite{SozioGionisKDD2010,SunPVLDB2025GICS}; it emphasizes efficient per-query answers (often under constraints such as size/connectivity), whereas we focus on global $(r,s)$-nucleus peeling with efficient dynamic index updates.
Finally, clique enumeration/counting are key primitives for higher-order cohesiveness; recent work revisits scalable $k$-clique counting by combining exact counting with sampling~\cite{YeSIGMOD2022Lightning,chen2025covering,yang2021huge,lai2019distributed}.
}

\section{Conclusion and future work} \label{sec:conclusion}

In this paper, we revisited nucleus decomposition through two complementary layers.
%
First, we introduced the \textit{Clique Path Index} (\cpi), a counting-based representation that eliminates explicit $s$-clique enumeration.
%
Second, we developed \emph{Differential Path Contribution Maintenance} (DPCM).
%
It updates exact supports by maintaining path-local sufficient states and differential contributions.
%
It does not reconstruct affected clique neighborhoods.
%
Together, these two layers yield exact regime-specific instantiations for $r{=}1$, $r{=}2$, and $r{\ge}3$.
%
They also preserve efficient hierarchy construction.
%
Extensive experiments show clear speedups for both decomposition and hierarchy construction.
%
They also show that our method scales to larger $s$.

\myaddRA{
Our counting-and-indexing viewpoint together with DPCM suggests natural future directions.
%
First, it is amenable to parallelization.
%
\textsc{BuildIndex} and \textsc{InitialCounting} can parallelize over pivot vertices or seed $r$-cliques with a final reduction.
%
\textsc{UpdateIndex} and \textsc{UpdateSupport} can process removed $r$-cliques in batches by using thread-local buffers and aggregated atomic updates.
%
More importantly, the DPCM formalism suggests a broader design space.
%
This design space includes new sufficient states, symbolic update operators, and regime-specific exact maintenance rules beyond the three instantiations studied here.
}
\myaddRC{
More broadly, CPI indexes the containment between small patterns and their supporting large patterns.
%
This lets us maintain supports by counting instead of explicit enumeration.
%
This idea extends beyond cliques~\cite{diestel2000graph}.
%
Possible targets include bicliques in bipartite graphs (related to $(\alpha,\beta)$-core)~\cite{tian2024abcore}, hyperedge-induced group patterns in hypergraphs~\cite{HyperCD,luo2024hierarchical,luo2025efficient}, and temporal $\Delta$-cliques~\cite{viard2016deltaclique,chen2024compressing}. 
%
We leave optimized multi-threaded implementations, richer DPCM state systems, and these extensions as future work.
}


\section{Acknowledgements}

% This work was partially supported by the NSFC Project 62302417.
Dong Wen is supported by ARC DP230101445 and ARC DE240100668.
%
Wenjie Zhang is supported by ARC DP230101445 and ARC FT210100303.
%
Xuemin Lin is supported by NSFC U2241211, NSFC U20B2046, and GuangDong Basic and Applied Basic Research Foundation 2019B1515120048.
